% Corrected French for printed p.158.
% Corrections: dus; abélienne; advance the duplicated resolution indices.
\SourcePage{158}

\LayerTarget{tohoku:construction:inverse-image-comparison:p2}
\noindent se prolonge de façon unique en un homomorphisme de foncteurs
cohomologiques
\[
R\Gamma^Y_\Psi\longrightarrow (R\Gamma^X_\Phi)g.
\]
On a donc prouvé :

\LayerTarget{tohoku:prop:3.2.2}
\textsc{Proposition 3.2.2.} --- \textit{Soit $f$ une application continue
d'un espace $X$ dans un espace $Y$, $\Phi$ (resp. $\Psi$) un antifiltre de
parties fermées de $X$ (resp. $Y$), tels que $B\in\Psi$ implique
$f^{-1}(B)\in\Phi$. On peut de façon unique trouver, pour tout faisceau de
groupes abéliens $F$ sur $Y$, des homomorphismes}
\[
H^p_\Psi(Y,F)\longrightarrow H^p_\Phi(X,f^{-1}(F))
\qquad(-\infty<p<+\infty)
\]
\textit{de façon à obtenir un homomorphisme de foncteurs cohomologiques, se
réduisant à l'homomorphisme naturel en dimension 0.}

Ces homomorphismes seront appelés homomorphismes naturels. De leur unicité
résulte une propriété de transitivité évidente, dont l'énoncé est laissé au
lecteur.

\LayerTarget{tohoku:construction:restriction-homomorphisms}
En particulier, si $Y$ est une \textit{partie de $X$}, et si on pose
$H^p_\Psi(Y,F)=H^p_\Psi(Y,F\vert Y)$, où $F\vert Y$ désigne la
\og restriction\fg{} du faisceau $F$ à $Y$, on a des \og homomorphismes de
restriction\fg{}
\[
H^p_\Phi(X,F)\longrightarrow H^p_{\Phi\cap Y}(Y,F),
\]
où $\Phi\cap Y$ désigne la \textit{trace} de $\Phi$ sur $Y$.

\LayerTarget{tohoku:section:3.3}
\textbf{3.3. Critères d'acyclicité.} Les développements de ce numéro, très
commodes pour la suite, sont dus à Godement, et seront traités de façon
détaillée dans le livre de Godement [9] cité dans l'Introduction.

\LayerTarget{tohoku:lemma:3.3.1}
\textsc{Lemme 3.3.1.} --- \textit{Soit $F$ un foncteur covariant d'une
catégorie abélienne $C$ dans une autre $C'$, on suppose que tout objet de $C$
est isomorphe à un sous-truc d'un objet injectif. Soit $\mathbf M$ une classe
d'objets de $C$, satisfaisant aux conditions suivantes : (i) Pour tout
$A\in C$, il existe un monomorphisme de $A$ dans un $M\in\mathbf M$ (ii) Tout
$A\in C$, isomorphe à un facteur direct d'un $M\in\mathbf M$, appartient à
$\mathbf M$ (iii) Pour toute suite exacte
$0\longrightarrow M'\longrightarrow M\longrightarrow M''\longrightarrow0$
dans $C$ avec $M'$ et $M$ dans $\mathbf M$, $M''$ appartient aussi à
$\mathbf M$, et la suite
$0\longrightarrow F(M')\longrightarrow F(M)\longrightarrow
F(M'')\longrightarrow0$ est aussi exacte. Sous ces conditions, tout objet
injectif de $C$ est dans $\mathbf M$, et pour tout $M\in\mathbf M$, on a
$R^pF(M)=0$ pour $p>0$.}

Soit d'abord $I$ un objet injectif dans $C$, $I$ se plonge dans un
$M\in\mathbf M$ en vertu de (i), et est donc isomorphe à un facteur direct de
$M$ (puisque $I$ est injectif), donc est dans $\mathbf M$ en vertu de (ii).
Soit $M\in\mathbf M$, prouvons $R^pF(M)=0$ pour $p>0$ ; pour ceci considérons
une résolution droite de $M$ par des objets injectifs $C^i$ ($i\geq0$), soit
$Z^i$ le sous-truc des cycles dans $C^i$, il suffit de prouver que les suites
\[
0\longrightarrow F(M)\longrightarrow F(C^0)\longrightarrow F(Z^1)
\longrightarrow0,\qquad
0\longrightarrow F(Z^1)\longrightarrow F(C^1)\longrightarrow F(Z^2)
\longrightarrow0,\quad\text{etc.}
\]
sont exactes, et pour ceci il suffit de prouver en vertu de (iii) que les
$Z^i$ et $C^i$ sont dans $\mathbf M$. On le sait déjà pour les $C^i$
puisqu'ils sont injectifs, et pour les $Z^i$ cela résulte de (iii) par
récurrence sur $i$.

\LayerTarget{tohoku:cor:3.3.1}
\textsc{Corollaire.} --- \textit{Pour tout $A\in C$, on peut calculer les
$R^pF(A)$ à l'aide d'une résolution (quelconque) de $A$ par des
$C^i\in\mathbf M$.}

En effet, une telle résolution existe en vertu de (i), et elle permet de
calculer les $R^pF(A)$ puisque les $C^i$ sont $F$-acycliques en vertu du
lemme.

\TohokuContinuation{159}

